\documentclass[aspectratio=169]{beamer}

% --- Toggle para version impresa (PDF) vs presentacion en vivo ---
% Cambiar \pptprinttrue por \pptprintfalse para version con animaciones embebidas
\newif\ifpptprint
\pptprinttrue

% --- Tema: solo colores, headline y footer custom ---
\usecolortheme{dolphin}
\setbeamertemplate{navigation symbols}{}
% --- Footline: solo el numero de diapositiva (guia 1.2) ---
% Las condiciones de simulacion NO van aca: estan una sola vez, en
% "Parametros del estudio". En cada diapositiva de figura solo se repite el
% parametro libre, a la derecha del titulo (ver \libre).
\setbeamertemplate{footline}{%
  \hbox to\paperwidth{\hfill
    \small\color{black!65}\insertframenumber\,/\,\inserttotalframenumber
    \hskip1.4em}%
  \vskip6pt
}

% --- Paquetes ---
\usepackage[utf8]{inputenc}
\usepackage[T1]{fontenc}
\usepackage[spanish,es-nodecimaldot]{babel}
\usepackage{amsmath,amssymb}
\usepackage{graphicx}
\usepackage{booktabs}
\usepackage{tikz}
\usepackage{hyperref}

% ============================================================
% Headline: grupos de puntos por seccion / estudio
% ============================================================
% Un punto por diapositiva. Los rangos de abajo deben coincidir con los
% numeros de frame reales: si se agregan o quitan diapositivas hay que
% actualizar los pares (primero,ultimo) de cada \dgrp.
\newcommand{\dgrp}[3]{% #1 = primer frame, #2 = ultimo frame, #3 = etiqueta
  \pgfmathtruncatemacro{\nd}{#2-#1+1}%
  \foreach \fr [count=\k from 0] in {#1,...,#2} {%
    \pgfmathsetmacro{\x}{\bx + \k*\dgap}%
    \ifnum\fr=\cf
      \fill[white] (\x,0) circle(\rad);
      \draw[white,line width=0.8pt] (\x,0) circle(\rad+1.2pt);
    \else\ifnum\fr<\cf
      \fill[white] (\x,0) circle(\rad);
    \else
      \draw[white,line width=0.4pt] (\x,0) circle(\rad);
    \fi\fi
  }%
  \node[white,font=\tiny\bfseries]
    at ({\bx+(\nd-1)*\dgap/2},-0.34) {#3};%
  \pgfmathsetmacro{\bxtmp}{\bx + (\nd-1)*\dgap + \sgap}%
  \let\bx\bxtmp
}

\setbeamertemplate{headline}{%
  \begin{beamercolorbox}[ht=2.9em,dp=0.1em,leftskip=0pt,rightskip=0pt]{section in head/foot}%
    \centering
    \begin{tikzpicture}
      \pgfmathtruncatemacro{\cf}{\insertframenumber}
      \def\rad{1.5pt}
      \def\dgap{0.19}
      \def\sgap{0.55}
      \pgfmathsetmacro{\bx}{0}

      % OJO: estos rangos estan hardcodeados. Si se agregan o quitan
      % diapositivas hay que actualizarlos.
      \dgrp{2}{5}{Intro.}
      \dgrp{6}{7}{Impl.}
      \dgrp{8}{12}{Sim.}
      \dgrp{13}{32}{Resultados}
      \dgrp{33}{34}{Conc.}
    \end{tikzpicture}
  \end{beamercolorbox}
}

% ============================================================
% Comandos de notacion y de armado de diapositivas
% ============================================================
\newcommand{\vv}[1]{\mathbf{#1}}      % vectores en negrita y sin italica
\newcommand{\uu}[1]{\mathrm{#1}}       % unidades en romano

% Diapositiva separadora de seccion (solo el titulo, sin numerar la seccion)
\newcommand{\seccion}[1]{%
  \begin{frame}
    \vfill
    \begin{center}{\Huge\bfseries #1}\end{center}
    \vfill
  \end{frame}
}

% Parametro libre de la diapositiva: va a la derecha del titulo, en negro.
% Es lo unico que se repite por diapositiva; el resto de las condiciones esta
% una sola vez en "Parametros del estudio".
% El titulo de beamer va en \raggedright, o sea \rightskip = 0pt plus 1fil.
% Un \hfill comun (1fil) se reparte el sobrante a medias con ese rightskip y el
% parametro queda a media asta; con 1filll gana y se pega al margen derecho.
\newcommand{\libre}[1]{\hskip 0pt plus 1filll\relax{\small\normalfont\color{black}#1}}

% Figura unica, al maximo tamano posible.
\newcommand{\figuno}[1]{%
  \vspace{-0.6em}
  \begin{center}
    \includegraphics[width=\textwidth,height=0.84\textheight,keepaspectratio]{#1}
  \end{center}
}

% Dos figuras lado a lado, cada una con su titulo justo debajo. El par se centra
% en \paperwidth (no en \textwidth) para ganar ancho; el \hss de los costados
% absorbe el desborde sin que LaTeX avise por overfull.
\newcommand{\figdos}[4]{% #1 = izq, #2 = titulo izq, #3 = der, #4 = titulo der
  \begin{center}
    \hbox to \textwidth{\hss
      \begin{minipage}[t]{0.492\paperwidth}
        \centering
        \includegraphics[width=\linewidth,height=0.78\textheight,keepaspectratio,
          trim=15 10 45 25,clip]{#1}
        \par\vspace{0.3em}{\small #2}
      \end{minipage}%
      \hspace{0.008\paperwidth}%
      \begin{minipage}[t]{0.492\paperwidth}
        \centering
        \includegraphics[width=\linewidth,height=0.78\textheight,keepaspectratio,
          trim=15 10 45 25,clip]{#3}
        \par\vspace{0.3em}{\small #4}
      \end{minipage}%
    \hss}
  \end{center}
}

% Placeholder para figuras todavia no generadas: caja con la ruta esperada.
\newcommand{\figplaceholder}[2]{%
  \begin{center}
    \fcolorbox{black!40}{black!10}{%
      \parbox[c][#1][c]{0.92\textwidth}{%
        \centering\small\textcolor{black!50}{FIGURA PENDIENTE\\\ttfamily\detokenize{#2}}}}
  \end{center}
}

% Dos animaciones por diapositiva. Cada columna deja el parametro arriba,
% el video centrado y el enlace explicito debajo (guia 2.4.8).
\newcommand{\animdos}[4]{% #1 = parametro izq, #2 = link izq,
                         % #3 = parametro der, #4 = link der
  \vspace{-0.3em}
  \begin{columns}[T,onlytextwidth]
    \column{0.485\textwidth}
      \centering
      {\large\bfseries #1}\par\vspace{0.5em}
      \fbox{\parbox[c][0.49\textheight][c]{0.92\linewidth}{%
        \centering\color{gray}
        \ifpptprint Fotograma del video\else Animación embebida\fi}}
      \par\vspace{0.5em}
      {\scriptsize Video: \href{#2}{#2}}

    \column{0.485\textwidth}
      \centering
      {\large\bfseries #3}\par\vspace{0.5em}
      \fbox{\parbox[c][0.49\textheight][c]{0.92\linewidth}{%
        \centering\color{gray}
        \ifpptprint Fotograma del video\else Animación embebida\fi}}
      \par\vspace{0.5em}
      {\scriptsize Video: \href{#4}{#4}}
  \end{columns}
}

% Una animacion por diapositiva: fotograma representativo y el enlace al
% video debajo (guia 2.4.8).
\newcommand{\animuno}[2]{% #1 = fotograma, #2 = link
  \vspace{-0.6em}
  \begin{center}
    \includegraphics[width=\textwidth,height=0.74\textheight,keepaspectratio]{#1}
    \par\vspace{0.3em}{\scriptsize Video: \href{#2}{#2}}
  \end{center}
}

% --- Metadatos ---
\title{Billar-Metegol}
\subtitle{Simulación dirigida por eventos y búsqueda de la configuración óptima}
\author{\texorpdfstring{%
  Juan Bautista Albertoni Salini --- 64189\\[0.25em]
  Santiago Devesa --- 64223\\[0.25em]
  Junior Rambau --- 64461%
}{Juan Bautista Albertoni Salini 64189; Santiago Devesa 64223; Junior Rambau 64461}}
\institute{(72.25) Simulación de Sistemas --- ITBA}
\date{Septiembre 2026}

% ============================================================
\begin{document}

% --- Frame 1: Titulo ---
\begin{frame}
  \titlepage
\end{frame}

\seccion{Introducción}

\section{Introducción}

\begin{frame}{El sistema}
  \begin{itemize}
    \setlength{\itemsep}{1.2em}
    \item Mesa de metegol con un arco en cada pared corta
    \item Obstáculos fijos que guían a las pelotas hacia los arcos
    \item $N$ esferas rígidas idénticas (masa $m_i$, radio $R_i$)
    \item Sin fuerzas de largo alcance ni rozamiento
  \end{itemize}
\end{frame}

\begin{frame}{Modelo matemático: partículas}
  \small
  \setlength{\abovedisplayskip}{3pt}
  \setlength{\belowdisplayskip}{3pt}
  \setlength{\abovedisplayshortskip}{3pt}
  \setlength{\belowdisplayshortskip}{3pt}
  \begin{itemize}
    \setlength{\itemsep}{1em}
    \item Vuelo libre de las partículas:
      \[
        x_i(t) = x_i(0) + v_{x_i}\,t, \qquad
        y_i(t) = y_i(0) + v_{y_i}\,t
      \]
    \item Colisión elástica $i$--$j$
      ($\sigma = R_i + R_j$;
      $\Delta\mathbf{v} = \mathbf{v}_j - \mathbf{v}_i$,
      $\Delta\mathbf{r} = \mathbf{r}_j - \mathbf{r}_i$):
      \smallskip
      \[
        J = \frac{2\,m_i m_j\,
            (\Delta\mathbf{v} \cdot \Delta\mathbf{r})}
            {\sigma\,(m_i + m_j)}
      \]
    \item Actualización post-choque:
      \begin{align*}
        vx_i^{d} &= vx_i^{a} + J_x/m_i, \qquad &
        vx_j^{d} &= vx_j^{a} - J_x/m_j, \\
        vy_i^{d} &= vy_i^{a} + J_y/m_i, \qquad &
        vy_j^{d} &= vy_j^{a} - J_y/m_j.
      \end{align*}
  \end{itemize}
\end{frame}

\begin{frame}{Modelo matemático: contorno}
  \small
  \setlength{\abovedisplayskip}{3pt}
  \setlength{\belowdisplayskip}{3pt}
  \setlength{\abovedisplayshortskip}{3pt}
  \setlength{\belowdisplayshortskip}{3pt}
  \begin{itemize}
    \setlength{\itemsep}{1.2em}
    \item Pared u obstáculo fijo ($m \to \infty$), colisión elástica:
      \[
        vx_i^{d} = -vx_i^{a} \quad \text{(pared vertical)}, \qquad
        vy_i^{d} = -vy_i^{a} \quad \text{(pared horizontal)}.
      \]
    \item Obstáculo circular fijo ($m \to \infty$), colisión elástica
      ($c_n = c_t = 1$):
      \smallskip
      \[
        \mathbf{\hat{e}}_n = (\cos\alpha, \sin\alpha), \qquad
        \mathbf{\hat{e}}_t = (-\sin\alpha, \cos\alpha)
      \]
      \medskip
      \[
        \mathbf{v}^d = \mathbf{R}(-\alpha)\, \mathbf{S}(c_n, c_t)\,
            \mathbf{R}(\alpha)\, \mathbf{v}^a
      \]
      \medskip
      \[
        \mathbf{R}(\alpha)
          = \begin{pmatrix}
              \cos\alpha & \sin\alpha \\
              -\sin\alpha & \cos\alpha
            \end{pmatrix}, \qquad
        \mathbf{S}(c_n, c_t)
          = \begin{pmatrix} -c_n & 0 \\ 0 & c_t \end{pmatrix}
      \]
  \end{itemize}
\end{frame}

\seccion{Implementación}

\section{Implementación}

\begin{frame}{Arquitectura del código}
  \figuno{entrega/uml/modulos.png}
\end{frame}

\seccion{Simulaciones}

\section{Simulaciones}

\begin{frame}{Geometría del sistema}
  \begin{columns}[T,onlytextwidth]
    \column{0.55\textwidth}
      \begin{itemize}
        \setlength{\itemsep}{0.6em}
        \item Mesa billar--metegol: $L \times W = 1.20 \times 0.68$ m.
        \item Arco: largo $d = 0.20$ m en cada pared corta.
        \item $N$ partículas idénticas: radio $r = 0.0175$ m, masa
          $m = 0.025$ kg.
        \item Rapidez inicial $v_0 = 1$ m/s.
        \item Obstáculos circulares fijos: opcionales.
      \end{itemize}

    \column{0.45\textwidth}
      \hfill
      \begin{tikzpicture}[x=4.3cm,y=4.3cm]
        \fill[black!6] (0,0) rectangle (1.2,0.68);
        \draw[thick] (0,0) rectangle (1.2,0.68);
        \draw[red!70!black,line width=2.2pt]
          (0,0.24) -- (0,0.44)  (1.2,0.24) -- (1.2,0.44);
        \fill[black!60]
          (0.55,0.15) circle(0.05)
          (0.55,0.53) circle(0.05)
          (0.80,0.34) circle(0.05);
        \foreach \x/\y/\ang in {0.10/0.08/60, 0.18/0.58/150, 0.32/0.42/210, 0.40/0.15/20, 0.26/0.26/300, 0.66/0.12/140, 0.76/0.50/200, 0.98/0.34/10, 0.05/0.34/80}
          {%
            \fill[blue!70!black] (\x,\y) circle(2.4pt);
            \draw[->,blue!70!black] (\x,\y) -- ++(\ang:0.075);
          }
      \end{tikzpicture}
  \end{columns}
\end{frame}

\begin{frame}{Observables}
  \begin{itemize}
    \setlength{\itemsep}{0.9em}
    \item $N_g(t)$: partículas que ya atravesaron un arco.
    \item $F_u(t) = N_g(t)/N$: fracción de partículas usadas.
    \item $t_{90}$: primer instante con $F_u(t) \ge 0.9$.
    \item $\langle t_{90} \rangle_n$: promedio de $t_{90}$ sobre $n$ realizaciones.
    \item $\mathrm{DCM}(t)$: desplazamiento cuadrático medio
      \[
        \mathrm{DCM}(t) = \frac{1}{N}\sum_{i=1}^{N}
          \left|\mathbf{r}_i(t) - \mathbf{r}_i(0)\right|^2
          \;\approx\; 4\,D\,t,
      \]
      con $D$ el coeficiente de difusión, ajustado en el tramo lineal.
  \end{itemize}
\end{frame}

\begin{frame}{Configuraciones estudiadas}
  \begin{center}
    \begin{tabular}{@{}llll@{}}
      \toprule
      Configuración & Obstáculos & Método & Se barre \\
      \midrule
      Mesa vacía        & ninguno              & baseline          & --- \\
      Círculo           & 1, centrado          & barrido exhaustivo & radio $R$ \\
      Quad              & $K$ c. simetría quad & Algoritmo genético (GA) & $K \in \{5,9,13,17\}$ \\
      Pared             & pared de círculos    & Algoritmo genético (GA) & grano, puntos de perfil \\
      \bottomrule
    \end{tabular}
  \end{center}
\end{frame}

\begin{frame}{Algoritmo genético}
  \renewcommand{\arraystretch}{1.7}
  \begin{center}
    \begin{tabular}{@{}l l@{}}
      \textbf{Fitness} &
        $f(s) = \begin{cases}
          t_{90}(s) & F_u \ge 0.9 \\
          t_{\max} + \alpha\,\bigl(90 - N_g(t_{\max})\bigr) & \text{si no}
        \end{cases}$ \\
      & $\langle t_{90} \rangle_5 = \frac{1}{5}\sum_{j=1}^{5} f(s_j)$, \ 5 semillas por generación \\
      \textbf{Selección} & Torneo de 3 \\
      \textbf{Cruza} & BLX-$\alpha$: \ $g_h \sim U\bigl[g_{\min} - 0.3\,\Delta,\ g_{\max} + 0.3\,\Delta\bigr]$,
        \ $\Delta = g_{\max} - g_{\min}$ \\
      \textbf{Mutación} & $g' = g + \mathcal{N}(0, \sigma^2)$ con prob.~0.3 por gen,
        \ $\sigma$ decrece con las generaciones \\
      \textbf{Elitismo} & 2 mejores \\
      \textbf{Validación} & Ganadores: $\langle t_{90} \rangle_{100}$\\
    \end{tabular}
  \end{center}
\end{frame}

\seccion{Resultados}

\section{Resultados}

\begin{frame}{Animaciones}
  \animdos{Mesa vacía}{https://youtu.be/TODO}{Configuración óptima}{https://youtu.be/TODO}
\end{frame}

\begin{frame}{Tiempo de ejecución vs Cantidad de partículas\libre{Sin obstáculos, $t_f = 30$ s}}
  \figuno{entrega/1.1/tiempo_ejecucion_vs_n_hex.png}
\end{frame}

% --- 1.2: circulo -> quad -> pared -> comparacion ---

\begin{frame}{Círculo: codificación\libre{$R \in \{0.05, 0.10, 0.15, 0.20, 0.25, 0.30, 0.34\}$ m}}
  \figuno{entrega/1.2/codificacion_circulo.png}
\end{frame}

\begin{frame}{$\langle t_{90} \rangle_{100}$ vs Radio del círculo}
  \figuno{entrega/1.2/circulo_t90_vs_radio.png}
\end{frame}

\begin{frame}{Círculo: mejor resultado}
  \animuno{entrega/1.2/circulo_mejor_caso_frame.png}{https://youtu.be/TODO}
\end{frame}

\begin{frame}{Quad: codificación\libre{$K \in \{5, 9, 13, 17\}$, pobl.~48, 60 gen.}}
  \figuno{entrega/1.2/codificacion_quad.png}
\end{frame}

\begin{frame}{Quad: geometrías propuestas por el GA}
  \figuno{entrega/1.2/quad_geometrias.png}
\end{frame}

\begin{frame}{Quad: evolución del ganador\libre{$K = 9$}}
  \figuno{entrega/1.2/quad_snapshots_k9.png}
\end{frame}

\begin{frame}{Quad: mejor resultado}
  \animuno{entrega/1.2/quad_k9_mejor_caso_frame.png}{https://youtu.be/TODO}
\end{frame}

\begin{frame}{Pared: codificación}
  \figuno{entrega/1.2/codificacion_pared.png}
\end{frame}

\begin{frame}{Pared: hiperparámetros\libre{Pobl.~40, 80 gen.}}
  \figuno{entrega/1.2/pared_hparams.png}
\end{frame}

\begin{frame}{Pared: geometrías por puntos de perfil\libre{Grano 0.018 m, pobl.~40, 80 gen.}}
  \figuno{entrega/1.2/pared_geometrias.png}
\end{frame}

\begin{frame}{Pared: evolución del ganador\libre{6 puntos, pobl.~64, 150 gen.}}
  \figuno{entrega/1.2/pared_snapshots.png}
\end{frame}

\begin{frame}{Pared: mejor resultado}
  \animuno{entrega/1.2/pared_mejor_caso_frame.png}{https://youtu.be/TODO}
\end{frame}

\begin{frame}{Comparación final}
  \figuno{entrega/1.2/comparacion_final.png}
\end{frame}

\begin{frame}{Fracción de usadas}
  \figplaceholder{0.6\textheight}{entrega/1.2/fu_vs_tiempo.png}
\end{frame}

\begin{frame}{DCM vs Tiempo}
  \figuno{entrega/1.3/dcm_comparacion.png}
\end{frame}

\begin{frame}{$\langle D \rangle$ vs $\langle t_{90} \rangle$}
  \figuno{entrega/1.3/t90_vs_difusion.png}
\end{frame}

\begin{frame}{Configuración elegida\libre{Competencia}}
  % TODO: esquema de la configuración final y su $\langle t_{90}\rangle$.
\end{frame}

\seccion{Conclusiones}

\section{Conclusiones}

\begin{frame}{Conclusiones}
  \begin{itemize}
    \setlength{\itemsep}{1em}
    \item % TODO
    \item % TODO
    \item % TODO
  \end{itemize}
\end{frame}

\begin{frame}
  \vfill
  \begin{center}{\Huge\bfseries ¡Gracias!}\end{center}
  \vfill
\end{frame}

\end{document}
